\documentclass[11pt,a4paper]{article}

\usepackage[utf8]{inputenc}
\usepackage[T1]{fontenc}
\usepackage[french]{babel}
\usepackage{geometry}
\usepackage{titlesec}
\usepackage{enumitem}
\usepackage{hyperref}

\geometry{margin=2.5cm}

\title{\textbf{Plateforme Agents IA -- Dashboard LLM \\ 
QALITAS et GMAO PRO}}
\author{TIM Techniques industrielles et Management}
\date{TBR. TIM.IA.S11/2025-2026}

\begin{document}

\maketitle

\section*{DOCUMENT STRICTEMENT CONFIDENTIEL}

\section{Système d’Agents IA pour le Monitoring \& la Gouvernance des Tokens LLM}

Système permettant de mesurer précisément la consommation de tokens LLM, maîtriser les coûts IA, prévenir les dérives, optimiser l’usage des prompts et modèles, assurer une facturation transparente aux clients.

\subsection*{Vue d’ensemble}

\textbf{Chaîne de valeur :}

\begin{center}
Mesurer $\rightarrow$ Calculer $\rightarrow$ Détecter $\rightarrow$ Expliquer $\rightarrow$ Optimiser $\rightarrow$ Anticiper $\rightarrow$ Gouverner $\rightarrow$ Piloter $\rightarrow$ Visualiser $\rightarrow$ Facturer
\end{center}

\textbf{Agents :}

\begin{itemize}
\item M1 --- Collecte \& Normalisation des Logs LLM
\item M2 --- Agrégation \& Calcul des Coûts Tokens
\item M3 --- Détection de Dérives \& Anomalies
\item M4 --- Explication \& Recommandations (LLM)
\item M5 --- Optimisation Prompts \& Modèles
\item M6 --- Prévision Budgétaire \& Capacity Planning
\item M7 --- Gouvernance \& Quotas IA
\item M8 --- Assistant Dashboard Tokens \& Coûts (LLM)
\item M9 --- Générateur \& Orchestrateur du Dashboard
\item M10 --- Facturation \& Reporting Client IA
\end{itemize}

\subsection*{Valeur globale}

\begin{itemize}
\item Maîtrise et monétisation des coûts IA
\item Gouvernance IA robuste (quotas, règles, audit)
\item Transparence totale client
\item Scalabilité sans dérive financière
\item Différenciation SaaS majeure
\end{itemize}

\textbf{Message clé :} Ce dispositif transforme l’IA d’un centre de coûts opaque en un actif mesuré, optimisé, gouverné et facturable.

\newpage

\section{AGENT M1 --- Collecte \& Normalisation des Logs LLM}

\subsection{Rôle stratégique}

L’agent M1 est le point de passage obligatoire de tous les appels LLM. Il transforme chaque appel IA en transaction mesurable, traçable et facturable.

\begin{itemize}
\item Sans M1 : consommation opaque
\item Avec M1 : consommation pilotable et monétisable
\end{itemize}

\subsection{Mission}

\begin{itemize}
\item Centraliser tous les appels LLM
\item Mesurer la consommation de tokens
\item Normaliser les logs
\item Alimenter le dashboard
\item Garantir la gouvernance et l’auditabilité
\end{itemize}

\subsection{Architecture}

\begin{center}
Modules $\rightarrow$ M1 (LLM Gateway) $\rightarrow$ Provider LLM
\end{center}

\textbf{Principe :} Aucun module n’appelle directement un LLM.

\subsection{Workflow}

\begin{enumerate}
\item Interception (trace, métadonnées)
\item Pré-contrôle (quotas, règles)
\item Appel LLM
\item Comptabilisation des tokens
\item Normalisation des données
\item Sécurité \& conformité
\item Persistance \& événements
\end{enumerate}

\subsection{KPI}

\begin{itemize}
\item Nombre d’appels
\item Tokens totaux
\item Coût IA
\item Latence
\item Taux d’erreur
\end{itemize}

\newpage

\section{AGENT M2 --- Agrégation \& Calcul des Coûts Tokens}

\subsection{Rôle stratégique}

M2 transforme les logs techniques en indicateurs financiers exploitables.

\begin{itemize}
\item M1 mesure
\item M2 calcule
\end{itemize}

\subsection{Mission}

\begin{itemize}
\item Consolider les logs
\item Calculer les coûts
\item Agréger par client, agent, modèle
\item Produire des tables BI
\end{itemize}

\subsection{Formules de coût}

\begin{itemize}
\item coût input = tokens\_input $\times$ prix\_input
\item coût output = tokens\_output $\times$ prix\_output
\item coût total = input + output
\end{itemize}

\subsection{KPI}

\begin{itemize}
\item Coût total IA
\item Coût par client
\item Coût par agent
\item Coût par modèle
\end{itemize}

\newpage

\section{AGENT M3 --- Détection de Dérives}

\subsection{Rôle}

M3 détecte les anomalies de consommation.

\begin{itemize}
\item M1 mesure
\item M2 calcule
\item M3 alerte
\end{itemize}

\subsection{Types de dérives}

\begin{itemize}
\item Volumétriques
\item Par appel
\item Par agent
\item Par modèle
\item Financières
\end{itemize}

\subsection{Méthodes}

\begin{itemize}
\item Règles métiers
\item Z-score
\item Comparaison temporelle
\end{itemize}

\subsection{Outputs}

\begin{itemize}
\item Alertes
\item Table anomalies
\item Score de criticité
\end{itemize}

\end{document}